\begin{frontmatter}
\title{Response to the Referee: Neural Bellman Operators}
\runtitle{NBO Response}
\end{frontmatter}
\section*{Revision and status}
We retain Neural Bellman Operators as the subject of the paper and pursue the
NBO contribution requested in the report. The revision brings the current
algorithm, general theory, capital-model analysis, and empirical comparison into
one article. The original economic applications and their complete derivations
remain in the supplement. The exact preceding numerical record is preserved,
including every unfavorable direct and mechanism comparison.

This response is a working template. The mathematical and numerical responses
below remain prospective until the frozen experiment, its independent inference,
and the final manuscript references have been inserted. Editorial completion is
not reported as resolution of an empirical objection.

\section*{B1: The distinctive NBO component is not supported by the executed evidence}
\paragraph*{Response.} The concern concerns the learned evaluation block, not the validity of the common verifier. We retain Neural Bellman Operators as the paper's subject and preserve the unfavorable R14 comparisons. The new study must identify the contribution of continuation evaluation through direct economic and work comparisons; the current editorial restructuring is not evidence that this scientific requirement has been met.
\paragraph*{Evidence required for the final response.} Source-frozen training protocol and complete run ledger; Direct NBO-minus-Raw/DPO/HJB endpoints; Occupation mechanism and full-work tables.
\paragraph*{Planned locations.} \path{sec:method}, \path{sec:r15mechanism}, \path{sec:r15evidence}.

\section*{B2: The costate-to-welfare mechanism is not operationally closed in the experiment}
\paragraph*{Response.} We agree that training-design costate diagnostics do not instantiate a candidate-occupation welfare bound. The revision will state the exact continuation target, occupation law, action-gap account and payoff transfer used by the new experiment. Finite-grid derivatives will not be identified with continuous-time costates, and an unbounded bias term will not be recorded as zero.
\paragraph*{Evidence required for the final response.} Reviewed mechanism theorem and complete proof; Occupation samples with policy identity; Enclosed actor gaps and all bias/transfer terms; Computed payoff loss decomposition.
\paragraph*{Planned locations.} \path{sec:r15mechanism}, \path{app:r15proofs}, \path{sec:r15evidence}.

\section*{B3: The high-dimensional result is safe improvement, not an accurate solution of the control problem}
\paragraph*{Response.} The existing high-dimensional regret bound is an anchor-based bound and does not become a tight HJB solution certificate because the schedule gain is positive. The revision will display its magnitude together with the gain, and assess any new accuracy result against a separately specified reference or bounding account. Scalar policy losses retain their scalar finite-grid scope.
\paragraph*{Evidence required for the final response.} New bound/reference or approximation study, if obtained; Anchor-gap and gain decomposition; Scalar five-state diagnostic and domain/mesh record.
\paragraph*{Planned locations.} \path{sec:capital}, \path{sec:r15mechanism}, \path{sec:r15evidence}.

\section*{B4: The statistical statements are conditional on selected fitted objects, not statements about the methods' training distributions}
\paragraph*{Response.} R14 certifies selected fitted policies conditional on their training records. Those intervals do not establish expected method performance or equivalence. The new design will specify its training-randomness population and economic comparison margin before final evaluation and report that inference separately from the conditional policy certificates.
\paragraph*{Evidence required for the final response.} Frozen estimand and seed-sampling plan; Complete training replication ledger; Method-level uncertainty derivation and output.
\paragraph*{Planned locations.} \path{sec:r15design}, \path{sec:r15evidence}.

\section*{B5: The work comparison is not yet a complete matched-accuracy comparison}
\paragraph*{Response.} Equal simulator visits and retrospective checkpoint times are not complete matched-accuracy costs. The revision will retain those original measurements under their original labels and use an inclusive clock for the new stopping experiment. Any efficiency claim must compare the same economic target and include unsuccessful checks and all charged computation.
\paragraph*{Evidence required for the final response.} End-to-end clocks and timing boundary specification; Simulator transitions, derivative evaluations and update counts; Memory and failed-check ledger; Accuracy/work frontier with censored failures.
\paragraph*{Planned locations.} \path{sec:r15design}, \path{sec:r15work}.

\section*{B6: The benchmark set is too narrow to establish a general numerical-method contribution}
\paragraph*{Response.} The baseline set must isolate what the separate critic and actor add. The revision will specify the shared economic and computational contract and include a direct HJB or approximate-policy-iteration comparator and a feasible classical comparison. Additional named methods will be included only when they address the same problem and information set; exclusions require a precise reason.
\paragraph*{Evidence required for the final response.} Comparator implementations and common contract; Classical training-method comparison; Documented mathematical incompatibility for each excluded method.
\paragraph*{Planned locations.} \path{sec:method}, \path{sec:r15design}, \path{sec:r15evidence}.

\section*{B7: The manuscript's scope and economic payoff are not aligned with the result established by R14}
\paragraph*{Response.} The original economic applications remain part of the NBO research program. They will be presented through the same evaluation-improvement structure with application-specific assumptions and deviation criteria, rather than as interchangeable evidence of platform-wide computational superiority. The current article will be separated from its historical development record, with complete models and proofs preserved.
\paragraph*{Evidence required for the final response.} Current application/proof inventory; Full old-label preservation map; Current article and supplement; Economic interpretation and units table.
\paragraph*{Planned locations.} \path{sec:introduction}, \path{sec:economic-scope}, \path{sec:r15design}.

\section*{M1: Choose the paper's center}
\paragraph*{Response.} We take the NBO-centered route. The title and original economic subject are retained, while the revision must establish what the continuation-learning block contributes. The method-agnostic validity of the verifier is retained as a supporting result and is not substituted for this empirical obligation.
\paragraph*{Evidence required for the final response.} A single current algorithm, theorem, and evidence chain; Verified incremental-value finding, if obtained.
\paragraph*{Planned locations.} \path{sec:introduction}, \path{sec:method}, \path{sec:r15mechanism}, \path{sec:r15evidence}.

\section*{M2: Define a method-level estimand before additional experiments}
\paragraph*{Response.} The new protocol will identify its primary method-level estimand, training-stream population and economic margin before confirmatory evaluation. It will distinguish variation across fitted objects from independent final-path error and retain every failed replication.
\paragraph*{Evidence required for the final response.} Protocol commit before confirmatory data; Training and final-evaluation sampling contracts.
\paragraph*{Planned locations.} \path{sec:r15design}.

\section*{M3: Make the critic mechanism operational}
\paragraph*{Response.} The mechanism account will use the distribution and continuation target appearing in its theorem. Each regression, action and transfer term must have its own computed value or bound. The resulting economic term will be compared with direct policy differences for the same fitted objects; a diagnostic table alone will not be called a closed welfare mechanism.
\paragraph*{Evidence required for the final response.} Term-by-term numerical instantiation; Direct contrasts using matching policies and units.
\paragraph*{Planned locations.} \path{sec:r15mechanism}, \path{app:r15proofs}, \path{sec:r15evidence}.

\section*{M4: Report complete accuracy-versus-work frontiers}
\paragraph*{Response.} The new work frontier will be based on executed stopping decisions with all earlier unsuccessful checks charged. The original prefix-clock calculations remain labeled retrospective. Unattained targets will be retained as censored outcomes rather than assigned zero work.
\paragraph*{Evidence required for the final response.} Actual stopping trace; Complete cost schema and frontier rows.
\paragraph*{Planned locations.} \path{sec:r15work}.

\section*{M5: Add contribution-isolating baselines}
\paragraph*{Response.} The comparison design will isolate the separate learned actor and continuation approximation under a common economic model and final payoff evaluator. The classical low-dimensional method will enter as a policy generator, and any baseline that cannot obey the same contract will be accompanied by a specific compatibility explanation.
\paragraph*{Evidence required for the final response.} Implementation and hyperparameter manifest; Final payoff and work ledger for all included methods.
\paragraph*{Planned locations.} \path{sec:r15design}, \path{sec:r15evidence}.

\section*{M6: Separate the model-specific theorem from generic claims}
\paragraph*{Response.} The current theorem sequence will distinguish the general evaluation-improvement statements from the sharper capital result, whose log felicity, bounded actions, additive diffusion and spectral structure will be stated with the theorem. The abstract and contribution list will identify that model-specific scope.
\paragraph*{Evidence required for the final response.} Assumption-to-result table; Revised abstract/contribution paragraphs.
\paragraph*{Planned locations.} \path{sec:model}, \path{sec:theory}, \path{sec:capital}.

\section*{M7: Clarify the economic implementation claim}
\paragraph*{Response.} The implementation discussion will describe the actual measured-state controller, its internal memory and its randomization, separately from the exact-history representation theorem. The sensing allowance will be compared only with the gain of the same protected parent, grid and population. Neither controller will be described as ordinary memoryless state feedback.
\paragraph*{Evidence required for the final response.} Controller information contract; Protected-parent identity and mesh; Sensing transfer and gain-ratio table.
\paragraph*{Planned locations.} \path{sec:r15implementation}, \path{app:r15observation}, \path{sec:r15evidence}.

\section*{M8: Reduce and reorganize the presentation}
\paragraph*{Response.} The revision will be organized by economic and mathematical subject rather than by R6-R14 development stages. Each substantive application and proof will have a current location, while repeated development narrative and complete historical records remain in the pinned archive. A label- and source-level map will make every relocation inspectable.
\paragraph*{Evidence required for the final response.} EDITORIAL\_MAP.json; Pinned R14 root/include inventory; Current PDFs and visual QA.
\paragraph*{Planned locations.} \path{sec:method}, \path{sec:theory}, \path{sec:r15design}, \path{app:r15index}.

\section*{M9: Clean up method labels and evidence semantics}
\paragraph*{Response.} The original raw records will remain byte-preserved. A documented adapter will retain their source method field while classifying fixed\_raw as raw\_costate, and all new records will use distinct method identifiers. Schedule-relative improvement and direct method superiority will be separate typed fields and captions.
\paragraph*{Evidence required for the final response.} Schema with source\_method\_id preserved; Adapter regression checks; New publication row ledger with explicit estimand.
\paragraph*{Planned locations.} \path{sec:r15design}, \path{app:r15index}.
